% !PRINTED 276
\weberpage{276}
%<<CONT
a change of sign is likewise lost; from this we conclude the theorem:

\gesperrt{Between two successive roots of $P_m=0$ there lies one and only one root of $P_{m-1}=0$.}

\begin{center}
\S.\,88.

\gesperrt{Second example.}
\end{center}

As a second example we consider an equation which, for brevity, we shall call the \gesperrt{secular equation}, because the investigation of the secular perturbations of the planets first led to it\footnote{\gesperrt{Laplace}, Histoire de l'Acad\'emie des Sciences 1772. Among more recent works one may consult on this \gesperrt{Dziobek}, Theorie der Planetenbewegung, Leipzig 1888.}). It is indeed otherwise well known under this name, but it also occurs in many other investigations, e.g. in the determination of the principal axes of a surface of the second degree, and in the theory of small oscillations; its general analytic significance lies in the fact that it furnishes a particular transformation, the so-called orthogonal transformation, of a quadratic form into a sum of squares. We shall here take the equation as it stands, without reference to any application, and shall discuss it according to our general theorems.

Let $a_{i,k}$, where $i$ and $k$ run through the series of indices $1, 2, 3 \ldots n$, be any system of real quantities, and let

\weq{1}{a_{i,k} = a_{k,i}}

be assumed. We consider the symmetric determinant

\weq{2}{L_n(x) = \begin{vmatrix} a_{1,1}-x, & a_{1,2} & \ldots & a_{1,n} \\ a_{2,1}, & a_{2,2}-x & \ldots & a_{2,n} \\ \cdot & \cdot & \cdot & \cdot \\ a_{n,1}, & a_{n,2} & \ldots & a_{n,n}-x \end{vmatrix}}

which is an integral rational \gesperrt{function} of degree $n$ in $x$, in which $x^n$ has the coefficient $(-1)^n$. The equation $L_n(x)=0$ shall be the object of our consideration.

We form the chain of principal subdeterminants taken with alternating signs

\weq{3}{L_n(x), \; -L_{n-1}(x), \; L_{n-2}(x) \ldots, \; (-1)^{n-1}L_1(x), \; (-1)^n,}
%CONT>>
